\documentclass[10pt]{article}

\usepackage{amsmath,amssymb,amsthm}
\usepackage{url}

\newtheorem{definition}{Definition}
\newtheorem{problem}{Problem}
\newtheorem{theorem}{Theorem}
\newtheorem{remark}[theorem]{Remark}

\newcommand{\dE}{\delta_E}
\newcommand{\Occ}{\operatorname{Occ}^E}
\newcommand{\PILLAR}{\textnormal{PILLAR}}
\newcommand{\PMwE}{\textnormal{PMwith\hspace{0pt}Edits}}
\newcommand{\NPM}{\textnormal{New\hspace{0pt}Periodic\hspace{0pt}Matches}}
\newcommand{\DPM}{\textnormal{Dynamic\hspace{0pt}Puzzle\hspace{0pt}Matching}}
\newcommand{\val}{\operatorname{val}}
\newcommand{\Ot}{\widetilde{O}}
\newcommand{\per}{\operatorname{per}}

\hyphenation{Mo-zes}

\title{The Cole--Hariharan Conjecture - Hard}
\author{Nikita Gaevoy}
\date{}

\begin{document}

\maketitle

\section{Preliminaries}

Everything in this section is a summary of published work. Subsections
\ref{sec:strings}--\ref{sec:pillar} follow Charalampopoulos, Kociumaka and
Wellnitz~\cite{CKW20,CKW22}; Subsections
\ref{sec:npm}--\ref{sec:seaweed} follow~\cite{CKW22} and, through it,
Tiskin's framework~\cite{Tiskin08,Tiskin15,TiskinArXiv}. Nothing in it is new.

\subsection{Strings and edit distance}\label{sec:strings}

We work with strings over an arbitrary alphabet $\Sigma$. A string $X$ has
length $|X|$ and consists of the characters $X[0]$, $X[1]$, \ldots,
$X[|X|-1]$; for $0\le i\le j\le|X|$ we write
\[
X[\,i\mathinner{.\,.}j\,) := X[i]X[i+1]\cdots X[j-1]
\]
for its \emph{fragment} of length $j-i$. A string $S$ is \emph{primitive} if
$S=U^y$ holds for no string $U$ and no integer $y>1$; $\per(S)$ denotes the
shortest period of $S$; $S^\infty$ denotes the infinite repetition
$S\cdot S\cdots$.

\begin{definition}[Edit distance]
The \emph{edit distance} $\dE(X,Y)$ of two strings $X,Y$ is the minimum number
of insertions, deletions and substitutions of single characters required to
transform $X$ into $Y$.
\end{definition}

For strings $U,V$ we further write, following~\cite[Section 1.3]{CKW22},
\[
\begin{aligned}
\dE(U,V^*)&:=\min\bigl\{\dE\bigl(U,V^\infty[\,0\mathinner{.\,.}j\,)\bigr):j\in\mathbb{Z}_{\ge0}\bigr\},\\
\dE(U,{*}V^*)&:=\min\bigl\{\dE\bigl(U,V^\infty[\,i\mathinner{.\,.}j\,)\bigr):0\le i\le j\bigr\}
\end{aligned}
\]
for the minimum edit distance between $U$ and a prefix, respectively a
substring, of $V^\infty$.

\subsection{Pattern matching with edits}\label{sec:pmwe}

\begin{definition}[$k$-error occurrences, {\cite[Section 1.3]{CKW22}}]\label{def:occ}
Let $P$ be a string (the \emph{pattern}), let $T$ be a string (the \emph{text})
and let $k>0$ be an integer (the \emph{threshold}). The pattern $P$ has a
\emph{$k$-error occurrence} in $T$ at position $v$ if $\dE(P,T[\,v\mathinner{.\,.}w\,))\le k$
for some $w\ge v$. We write
\[
\Occ_k(P,T):=\bigl\{v:\exists_{w\ge v}\ \dE\bigl(P,T[\,v\mathinner{.\,.}w\,)\bigr)\le k\bigr\}
\]
for the set of starting positions of $k$-error occurrences of $P$ in $T$.
\end{definition}

\begin{definition}[\PMwE{}, {\cite[Section 1.3]{CKW22}}]\label{def:pmwe}
The \emph{pattern matching with edits} problem \PMwE$(P,T,k)$ takes as input a
pattern $P$ of length $m$, a text $T$ of length $n$, and a positive integer
$k\le m$, and asks for the set $\Occ_k(P,T)$.
\end{definition}

Throughout, $m$, $n$ and $k$ denote these three quantities, and $m\le n$.

\begin{remark}\label{rem:report}
A $k$-error occurrence is a \emph{starting position} $v$; the problem does not
ask for the corresponding endpoints $w$, nor for the alignments. Landau, Myers
and Schmidt~\cite{LMS98} show that all the fragments of $T$ at edit distance
at most $k$ from $P$ can be listed in $O(nk)$ time; see
Remark~\ref{rem:variants} for the status of this stronger output requirement in
the regime considered here.
\end{remark}

\begin{remark}\label{rem:output}
For the running times considered below, the output has to be represented
implicitly. Charalampopoulos, Kociumaka and Wellnitz~\cite[Main Theorem
7]{CKW20} show that $\Occ_k(P,T)$ always decomposes into $O(n/m\cdot k^3)$
disjoint arithmetic progressions with a common difference, and this is the
representation used by all the results quoted below. Unless $P$ is
approximately periodic, $|\Occ_k(P,T)|=O(n/m\cdot k^2)$ and the set may be
reported explicitly~\cite[Section 1]{CKW22}.
\end{remark}

\subsection{The \PILLAR{} model}\label{sec:pillar}

\begin{definition}[\PILLAR{} model, {\cite{CKW20}}, {\cite[Section 2.1]{CKW22}}]\label{def:pillar}
In the \PILLAR{} model one maintains a collection $\mathcal X$ of strings and
operates on fragments of the strings in $\mathcal X$, each represented by a
handle. The running time of an algorithm is the number of calls to the
following \emph{primitive \PILLAR{} operations}, plus any extra computation,
which is accounted for separately:
\begin{itemize}
    \item $\mathrm{Extract}(S,\ell,r)$: given a fragment $S$ and
    $0\le\ell\le r\le|S|$, return a handle to $S[\,\ell\mathinner{.\,.}r\,)$;
    \item $\mathrm{LCP}(S,T)$ and $\mathrm{LCP}^R(S,T)$: the length of the
    longest common prefix, respectively suffix, of $S$ and $T$;
    \item $\mathrm{IPM}(P,T)$: assuming $|T|\le2|P|$, the set of exact
    occurrences of $P$ in $T$, represented as an arithmetic progression with
    difference $\per(P)$;
    \item $\mathrm{Access}(S,i)$ and $\mathrm{Length}(S)$.
\end{itemize}
\end{definition}

In the standard setting, where $P$ and $T$ are given explicitly, each primitive
operation is implemented in $O(1)$ time after an $O(n)$-time
preprocessing~\cite[Section 7.1]{CKW22}; the same implementation serves the
internal setting of~\cite{KRRW15}. Implementations are also available in the
fully compressed setting, where the strings are given as straight-line
programs, and in the dynamic setting, where a collection of strings is
maintained under creation, splitting and concatenation~\cite[Sections 7.2,
7.3]{CKW22}. A \PILLAR{}-model bound therefore yields an algorithm in each of
these settings at once.

Of the derived operations of the \PILLAR{} toolbox of~\cite{CKW20} we shall
need only the following one.

\begin{theorem}[$\mathrm{Verify}(P,T,k,I)$, {\cite[Section 5]{CH02}},
{\cite[Fact 2.9]{CKW22}}]\label{thm:verify}
Let $P$ be a string of length $m$, let $T$ be a string, let $k\le m$ be a
positive integer and let $I$ be an interval of positions. Using
$O\bigl(k(k+|I|)\bigr)$ \PILLAR{} operations one can compute
\[
\bigl\{\bigl(\ell,\ \textstyle\min_{r}\dE(P,T[\,\ell\mathinner{.\,.}r\,))\bigr):
\ell\in\Occ_k(P,T)\cap I\bigr\}.
\]
\end{theorem}

\subsection{The periodic bottleneck}\label{sec:npm}

The structural analysis of~\cite{CKW20} reduces \PMwE{} to the case in which
both the pattern and the text are approximately periodic. That case is
isolated in~\cite{CKW22} as follows.

\begin{definition}[\NPM{}, {\cite[Section 1.3]{CKW22}}]\label{def:npm}
An instance of \NPM$(P,T,k,d,Q,\mathcal A_P,\mathcal A_T)$ consists of
\begin{itemize}
    \item a pattern $P$ of length $m$ and a threshold $k\in[\,0\mathinner{.\,.}m\,]$;
    \item a positive integer $d\ge 2k$;
    \item a text $T$ of length $n\in[\,m-k\mathinner{.\,.}3/2\,m+k\,)$;
    \item a primitive string $Q$ of length $q:=|Q|\le m/8d$;
    \item edit-distance alignments
    \[
    \mathcal A_P:P\rightsquigarrow Q^\infty[\,0\mathinner{.\,.}y_P\,)
    \qquad\text{and}\qquad
    \mathcal A_T:T\rightsquigarrow Q^\infty[\,x_T\mathinner{.\,.}y_T\,)
    \]
    with $x_T\in[\,0\mathinner{.\,.}q\,)$, of respective costs
    \[
    d_P:=\dE(P,{*}Q^*)=\dE(P,Q^*)\le d
    \qquad\text{and}\qquad
    d_T:=\dE(T,{*}Q^*)\le3d.
    \]
\end{itemize}
The output is the set $\Occ_k(P,T)$, represented as $O(d^3)$ disjoint
arithmetic progressions with difference $q$.
\end{definition}

\begin{theorem}[{\cite[Fact 1.1]{CKW22}}, after {\cite{CKW20}}]\label{thm:reduction}
Let $P$ be a pattern of length $m$, let $T$ be a text of length $n$ and let
$k\le m$ be a positive integer. One can compute a representation of
$\Occ_k(P,T)$ as $O(n/m\cdot k^3)$ disjoint arithmetic progressions with the
same difference in time $O(n/m\cdot k^3)$ in the \PILLAR{} model, plus the time
required to solve several instances
\[
\NPM(P_i,T_i,k_i,d_i,Q_i,\mathcal A_{P_i},\mathcal A_{T_i})
\]
with $\sum_i|P_i|=O(n)$ and, for each $i$, $|P_i|\le m$ and
$d_i=\lceil 8k/m\cdot|P_i|\rceil$.
\end{theorem}

\begin{remark}\label{rem:npmscaling}
Theorem~\ref{thm:reduction} is what makes \NPM{} the right subproblem to
optimise: since $d_i=\Theta(k/m\cdot|P_i|)$ and $\sum_i|P_i|=O(n)$, an
$O(d^{c})$-time algorithm for \NPM{} yields an $O(n/m\cdot k^{c})$-time
algorithm for \PMwE{} for every $c\ge1$, and polylogarithmic factors carry over
unchanged~\cite[proof of Main Theorem 2]{CKW22}.
\end{remark}

The heavy\slash light dichotomy of~\cite[Sections 5, 6]{CKW22} splits the positions
of $T$ according to a threshold $\eta$. Its definition rests on the
\emph{locked fragments} of~\cite[Lemma 6.9]{CKW20} and on a marking scheme
whose details~\cite[Definitions 5.5, 5.8]{CKW22} are not needed to state the
problems below; what is needed is the following interface.

\begin{theorem}[$\mathrm{Heavy}$, {\cite[Lemmas 5.9, 5.10]{CKW22}}]\label{thm:heavy}
Consider an instance of \NPM{} and an integer threshold $\eta\ge1$. The
associated set $H_\eta\subseteq[\,0\mathinner{.\,.}n-m+k\,]$ of \emph{heavy positions} can
be computed in $O(d^2\log\log d)$ time in the \PILLAR{} model and is
represented as the union of $O(d^2)$ disjoint integer ranges. Its size
satisfies
\[
|H_\eta|=O\bigl((d^2/\eta+d)(d+q)\bigr).
\]
The positions in $[\,0\mathinner{.\,.}n-m+k\,]\setminus H_\eta$ are called \emph{light}.
\end{theorem}

\subsection{Puzzles, seaweeds and dynamic puzzle matching}\label{sec:seaweed}

The material of this subsection is not needed to state the problem below; it
is used in the remarks that follow it, and it is the apparatus in which the
known approaches to the problem are expressed.

\begin{definition}[$\Delta$-puzzle, {\cite[Definition 1.4]{CKW22}}]\label{def:puzzle}
For $\Delta\in\mathbb{Z}_{\ge0}$, strings $S_1,\ldots,S_z$ with $z\ge2$ form a
\emph{$\Delta$-puzzle} if $|S_i|\ge\Delta$ for each $i\in[\,1\mathinner{.\,.}z\,]$ and
$S_i[\,|S_i|-\Delta\mathinner{.\,.}|S_i|\,)=S_{i+1}[\,0\mathinner{.\,.}\Delta\,)$ for each
$i\in[\,1\mathinner{.\,.}z\,)$. Its \emph{value} is
\[
\val_\Delta(S_1,\ldots,S_z):=S_1\cdot S_2[\,\Delta\mathinner{.\,.}|S_2|\,)\cdots
S_z[\,\Delta\mathinner{.\,.}|S_z|\,).
\]
\end{definition}

\begin{definition}[\DPM{}, {\cite[Section 1.3]{CKW22}}]\label{def:dpm}
An instance \DPM$(k,\Delta,\mathcal S_\beta,\mathcal S_\mu,\mathcal S_\varphi)$
is given by positive integers $k$ and $\Delta$ and by families
$\mathcal S_\beta$, $\mathcal S_\mu$, $\mathcal S_\varphi$ of \emph{leading},
\emph{internal} and \emph{trailing} pieces. The maintained object is a
sequence $\mathcal I=(U_1,V_1)\cdots(U_z,V_z)$ of pairs of strings such that
$U_1,V_1\in\mathcal S_\beta$, $U_z,V_z\in\mathcal S_\varphi$, and
$U_i,V_i\in\mathcal S_\mu$ for $i\in(\,1\mathinner{.\,.}z\,)$, and such that the
\emph{torsion} $\operatorname{tor}(\mathcal I):=\sum_{i=1}^{z}|U_i|-|V_i|$
satisfies $\operatorname{tor}(\mathcal I)\le\Delta/2-k$. The updates are the
deletion of a pair, the insertion of a pair, and the substitution of a pair.
A query returns
\[
\Occ_k(\mathcal I):=\Occ_k\bigl(\val_\Delta(U_1,\ldots,U_z),\
\val_\Delta(V_1,\ldots,V_z)\bigr)
\]
under the promise that $U_1,\ldots,U_z$ and $V_1,\ldots,V_z$ are
$\Delta$-puzzles.
\end{definition}

For a family $\mathcal S$ of strings, its \emph{median edit distance} is
\[
\dE(\mathcal S):=\min_{\hat S\in\Sigma^*}\sum_{S\in\mathcal S}\dE(S,\hat S).
\]

\begin{theorem}[{\cite[Theorem 1.5]{CKW22}}]\label{thm:dpm}
There is a data structure for the problem
\[
\DPM(k,\Delta,\mathcal S_\beta,\mathcal S_\mu,\mathcal S_\varphi)
\]
with $O(\Delta\log z\log\Delta)$-time updates and queries,
$O(\Delta z\log\Delta)$-time initialisation and
$O\bigl((d^3+\Delta^2d)\log^2(d+\Delta)\bigr)$-time preprocessing, where
$d=\dE(\mathcal S_\beta)+\dE(\mathcal S_\mu)+\dE(\mathcal S_\varphi)$ and $z$
is the length of the maintained sequence.
\end{theorem}

The data structure of Theorem~\ref{thm:dpm} is built on the seaweed monoid of
permutation matrices, also known as the monoid of sticky
braids~\cite{Tiskin08,Tiskin15,TiskinArXiv}. For strings $X,Y$ one considers
the alignment graph on $[\,0\mathinner{.\,.}|Y|\,]\times[\,0\mathinner{.\,.}|X|\,]$ in which horizontal
and vertical edges have weight $1$ and a diagonal edge has weight $0$ for a
match and $1$ for a substitution; then $\dE(X,Y[\,v\mathinner{.\,.}w\,))$ is the distance
from $(v,0)$ to $(w,|X|)$. All these distances are encoded in linear space by a
permutation matrix $P_{Y,X}$, the \emph{seaweed matrix}, and seaweed matrices
compose under concatenation of either argument by the \emph{seaweed product}
$\odot$, which is the implicit $(\min,+)$ product of the corresponding
unit-Monge matrices~\cite[Lemma 9.2]{CKW22}.

\begin{theorem}[{\cite[Theorem 8.2]{CKW22}}, after {\cite{Tiskin15,TiskinArXiv}}]\label{thm:steadyant}
For any two bounded permutation matrices $A,B$ of size $n$, the seaweed product
$A\odot B$ can be computed in $O(n\log n)$ time.
\end{theorem}

Since $P_{Y,X}$ encodes $\dE(X,Y)$, it cannot be computed in truly subquadratic
time. The contribution of~\cite[Section 8]{CKW22} is that, when only alignments
of cost at most $k$ matter, the relevant paths stay inside a diagonal band
$I=[\,-k\mathinner{.\,.}|Y|-|X|+k\,]$ of the alignment graph, and the seaweed matrix of the
graph restricted to that band admits an encoding of size $O(|I|)$; this defines
a \emph{restriction} operation $P\mapsto P|_I$ on bounded permutation matrices
that is computable in linear time and commutes with $\odot$~\cite[Lemmas 8.16,
8.20, 8.21]{CKW22}. The \DPM{} data structure expresses the restricted seaweed
matrix of $\val_\Delta(V_1,\ldots,V_z)$ against $\val_\Delta(U_1,\ldots,U_z)$
as a seaweed product of $z$ restricted matrices of size $O(\Delta)$, one per
pair, and maintains that product in a balanced binary tree; an update
recomputes $O(\log z)$ partial products by Theorem~\ref{thm:steadyant}, and a
query applies the SMAWK algorithm~\cite{AKMSW87} to the matrix at the
root~\cite[Lemma 9.5, Section 9]{CKW22}.

\section{Problem formulation}

\begin{problem}[The Cole--Hariharan conjecture]\label{prob:CH}
Prove that there is an algorithm which, given a text $T$ of length $n$, a
pattern $P$ of length $m$ and an integer threshold $k>0$, computes
$\Occ_k(P,T)$ in time
\[
\Ot\bigl(n+k^3\cdot n/m\bigr),
\]
that is, in time $O\bigl(n+k^3\log^{c}m\cdot n/m\bigr)$ for some constant
$c>0$. Better still, and the form in which the conjecture was stated by Cole
and Hariharan~\cite{CH02} and is quoted in~\cite[Section 1]{CKW22}, is the
sharper bound
\[
O\bigl(n+k^3\cdot n/m\bigr)
\]
without the hidden factors.
\end{problem}

\begin{remark}\label{rem:CHsoft}
The polylogarithmic form is asked for first because it already settles the
question the current techniques leave open: the state of the
art~\cite[Main Theorems 1 and 2]{CKW22} itself carries a subpolynomial factor
$\sqrt{\log m\log k}$. The sharp form is the stronger statement and implies it.
Both are open.
\end{remark}

\begin{remark}\label{rem:scope}
The scope stops at $k^3$, rather than at the conditional lower bound
$k^{2-\Omega(1)}$ that follows from SETH~\cite{BI18},~\cite[Section 1]{CKW22}, because $k^3$ is where the output
representation used throughout~\cite{CKW20,CKW22} runs out: that
representation writes $\Occ_k(P,T)$ using $\Ot(n/m\cdot k^3)$ bits
(Remark~\ref{rem:output}), and, as observed by Kociumaka, Nogler and
Wellnitz~\cite[Section 1]{KNW26}, it is therefore ``fundamentally incompatible
with reaching the lower bound''. Going below $k^3$ requires a different
encoding of the output; that such an encoding exists at all is the content
of~\cite{KNW24,KNW26}, which show that $O(n/m\cdot k\log(m|\Sigma|/k))$ bits
suffice, but no algorithm is known to exploit it.
\end{remark}

\begin{remark}\label{rem:variants}
Two variants are equally interesting and are not settled by the problems above.
First, the decision version, which asks only whether the set $\Occ_k(P,T)$ is
empty, and the approximate version, which allows reporting an arbitrary subset
of
$\Occ_{(1+\varepsilon)k}(P,T)\setminus\Occ_k(P,T)$, are proposed
in~\cite[Section 1.2]{CKW22} as possibly easier relaxations. Second, the
reporting version of Remark~\ref{rem:report}, which asks for all fragments of
$T$ at edit distance at most $k$ from $P$ rather than for their starting
positions only: the authors of~\cite{CKW22} write that the $O(k^4\cdot n/m)$
\PILLAR{} algorithm of~\cite{CKW20} can be generalised to report all such
fragments, whereas their own $\Ot(k^{3.5}\cdot n/m)$ solution ``does not seem to
generalise''~\cite[Section 1.2]{CKW22}. A solution for starting positions alone is a
solution to the problem as posed.
\end{remark}

\newcommand{\etalchar}[1]{$^{#1}$}
\begin{thebibliography}{KRRW15}

\bibitem[AKM{\etalchar{+}}87]{AKMSW87}
Alok Aggarwal, Maria~M. Klawe, Shlomo Moran, Peter~W. Shor, and Robert~E.
  Wilber.
\newblock Geometric applications of a matrix-searching algorithm.
\newblock {\em Algorithmica}, 2:195--208, 1987.

\bibitem[BI18]{BI18}
Arturs Backurs and Piotr Indyk.
\newblock Edit distance cannot be computed in strongly subquadratic time
  (unless {SETH} is false).
\newblock {\em SIAM Journal on Computing}, 47(3):1087--1097, 2018.

\bibitem[CH02]{CH02}
Richard Cole and Ramesh Hariharan.
\newblock Approximate string matching: A simpler faster algorithm.
\newblock {\em SIAM Journal on Computing}, 31(6):1761--1782, 2002.

\bibitem[CKW20]{CKW20}
Panagiotis Charalampopoulos, Tomasz Kociumaka, and Philip Wellnitz.
\newblock Faster approximate pattern matching: A unified approach.
\newblock In {\em 61st Annual IEEE Symposium on Foundations of Computer Science
  (FOCS 2020)}, pages 978--989, 2020.
\newblock Full version: arXiv:2004.08350.

\bibitem[CKW22]{CKW22}
Panagiotis Charalampopoulos, Tomasz Kociumaka, and Philip Wellnitz.
\newblock Faster pattern matching under edit distance: A reduction to dynamic
  puzzle matching and the seaweed monoid of permutation matrices.
\newblock In {\em 63rd Annual IEEE Symposium on Foundations of Computer Science
  (FOCS 2022)}, pages 698--707, 2022.
\newblock Full version: arXiv:2204.03087.

\bibitem[KNW24]{KNW24}
Tomasz Kociumaka, Jakob Nogler, and Philip Wellnitz.
\newblock On the communication complexity of approximate pattern matching.
\newblock In {\em 56th Annual ACM Symposium on Theory of Computing (STOC
  2024)}, 2024.
\newblock arXiv:2403.18812.

\bibitem[KNW26]{KNW26}
Tomasz Kociumaka, Jakob Nogler, and Philip Wellnitz.
\newblock The communication complexity of pattern matching with edits
  revisited.
\newblock In {\em 37th Annual Symposium on Combinatorial Pattern Matching (CPM
  2026)}, 2026.
\newblock arXiv:2604.15601.

\bibitem[KRRW15]{KRRW15}
Tomasz Kociumaka, Jakub Radoszewski, Wojciech Rytter, and Tomasz Wale{\'n}.
\newblock Internal pattern matching queries in a text and applications.
\newblock In {\em 26th Annual ACM-SIAM Symposium on Discrete Algorithms (SODA
  2015)}, pages 532--551, 2015.

\bibitem[LMS98]{LMS98}
Gad~M. Landau, Eugene~W. Myers, and Jeanette~P. Schmidt.
\newblock Incremental string comparison.
\newblock {\em SIAM Journal on Computing}, 27(2):557--582, 1998.

\bibitem[Tis08]{Tiskin08}
Alexander Tiskin.
\newblock Semi-local string comparison: Algorithmic techniques and
  applications.
\newblock {\em Mathematics in Computer Science}, 1(4):571--603, 2008.

\bibitem[Tis13]{TiskinArXiv}
Alexander Tiskin.
\newblock Semi-local string comparison: algorithmic techniques and
  applications, 2013.
\newblock arXiv:0707.3619v21.

\bibitem[Tis15]{Tiskin15}
Alexander Tiskin.
\newblock Fast distance multiplication of unit-{M}onge matrices.
\newblock {\em Algorithmica}, 71(4):859--888, 2015.

\end{thebibliography}

\end{document}
